\documentclass[../main.tex]{subfiles}

\begin{document}
\chapter{Lie Group and Lie Algebra}

\begin{definition}{}{}
    \begin{enumerate}
        \item (拓扑群)  设 \(  X  \)是一个拓扑空间.  若存在连续映射 \[
        \begin{aligned}
        \mu : X\times X&\to X\\ 
         \left( x,y \right)& \mapsto  xy  
        \end{aligned}
        \] 以及 \[
        \begin{aligned}
        i: X\to X,\quad x\mapsto x^{-1}  
        \end{aligned}
        \]满足
        \begin{enumerate}
            \item  \(  \left( xy \right)z= x\left( yz \right)    \).
            \item 存在 \(  e \in X  \) ,使得 \(  ex= xe= x    \) .
            \item \(  \forall x\in X, \exists y \in X  \),使得 \(  xy= yx= e  \) . 此时 \(  y  \)唯一, 记作 \(  x^{-1}   \) .   
        \end{enumerate}
    
    \item (李群) 若 \(  G  \)是 \(  C^{\infty}  \)流形. 其上有\(  C^{\infty}  \)映射 \[
    \mu : G\times G\to G,\quad i:G\to G
    \]   以及 \(  e \in G  \). 满足同1.的群的定义. 则称 \(  G  \)是一个李群.   
    \end{enumerate}
    
\end{definition}
\begin{remark}
    对于李群 \(G\).可以忘掉其微分结构, 将其视为一个拓扑群.它是第二可数,  Hausdorff且局部欧的. 
\end{remark}

\begin{definition}{}{}
    设 \(  G_1,G_2  \)是两个李群. 若 \(   \varphi :G_1\to G_2  \)既是群同态, 又是\(  C^{\infty}  \)映射. 则称\(   \varphi   \)是一个李群同态.     
\end{definition}
\begin{definition}{}{}
    对于 \(  a \in G  \) , 记   \[
    \begin{aligned}
     \ell_{a}: G&\to G\\ 
      x & \mapsto  ax
    \end{aligned}
    \]称为左平移. 记 \[
    \begin{aligned}
    r_{a}: G&\to G\\ 
     x\mapsto xa 
    \end{aligned}
    \]称为右平移. 
    记 \[
    \begin{aligned}
    i_{a}: G&\to G\\ 
     x&\mapsto  axa^{-1}  
    \end{aligned}
    \]称为共轭. 
\end{definition}

\begin{definition}{}{}
  设 \(  G,H  \)是两个李群. 若 \(  H  \)既是 \(  G  \)的子流形, 又是 \(  G  \)的子群. 则称 \(  H  \)为 \(  G  \)的一个李子群.         
\end{definition}


\begin{theorem}{}{}
    设 \(  G  \)是一个李群. \(  H  \)是 \(  G  \)的子群. 若 \(  H  \)是拟正则的嵌入子流形. 则它是李子群. 若 \(  H  \)还是正则嵌入的子流形. 则 \(  H  \)是一个闭的李子群. 
\end{theorem}
\begin{proof}
 考虑   \[
    \begin{aligned}
    \nu : G\times G&\to G\\ 
     \left( x,y \right)&\mapsto xy^{-1}   
    \end{aligned}
    \]是一个 \(  C^{\infty}  \)映射.  
    若考虑光滑映射 \(\tilde{\nu} =  \left. \nu \right|_{H\times H}: H\times H\to H\hookrightarrow G  \).  若 \(  j  \)是拟正则的嵌入映射,  取 \(  P = H\times H,u: H\times H\to H  \) . 则 \[
    \tilde{\nu }= j\circ u \text{ is } C^{\infty} \implies  u \text{ is } C^{\infty} \implies  H  \text{ is a Lie subgroup} 
    \]

    对于第二个断言. 由于正则子流形是局部闭的,  
\end{proof}
\end{document} 